\chapter{Rings and Arithmetic}\label{ch:b3:rings}

Ordinary integers factor uniquely into primes; so do polynomials
over a field. Are these two facts one theorem? This chapter answers
yes, and finds the exact hypotheses that make an ``arithmetic''
possible in a commutative ring: the chain
\[
\text{Euclidean} \;\Longrightarrow\; \text{principal}
\;\Longrightarrow\; \text{factorial (UFD)},
\]
with all implications proved and all converses refuted. The theory
is then tested where it earns its keep: the Gaussian integers
$\Z[\iu]$ (which will crack Fermat's two-squares theorem in the
weekend problem), polynomial rings in several variables (Gauss's
lemma, Eisenstein's criterion), and Noetherian rings, culminating
in Hilbert's basis theorem. Throughout, \emph{ring} means
commutative ring with unit $1 \neq 0$; the Year 2 volume's ideals
of $\Z$ and $K[X]$ are our two guiding examples.

\section{Ideals, quotients, and the isomorphism theorem}

\begin{definition}\label{def:b3:rings:ideal}
An \emph{ideal}\index{ideal} $I$ of a ring $A$ is an additive
subgroup such that $AI \subseteq I$. The \emph{quotient
ring}\index{quotient ring} $A/I$ is the quotient group $(A, +)/I$
with the multiplication $(a + I)(b + I) = ab + I$: well defined,
since changing $a$ to $a + x$, $b$ to $b + y$ ($x, y \in I$)
changes $ab$ by $ay + xb + xy \in I$. The projection $\pi \colon A
\to A/I$ is a surjective ring morphism with kernel $I$, and kernels
of ring morphisms are exactly the ideals.
\end{definition}

\begin{theorem}[First isomorphism theorem]\label{thm:b3:rings:firstiso}
If $f \colon A \to B$ is a ring morphism, then $\bar f\colon
A/\ker f \to \operatorname{im} f$, $a + \ker f \mapsto f(a)$, is a
ring isomorphism. More generally $f$ factors through $A/I$ for any
ideal $I \subseteq \ker f$. The ideals of $A/I$ are the $J/I$ for
$J \supseteq I$ an ideal of $A$ (correspondence theorem).
\end{theorem}

\begin{proof}
As for groups (\cref{thm:b3:groups:firstiso,%
thm:b3:groups:correspondence}), noting that all maps in sight also
respect products: $\bar f$ is well defined, bijective onto the
image, and multiplicative; the correspondence $J \mapsto J/I$,
$\bar J \mapsto \pi^{-1}(\bar J)$ preserves ideals in both
directions because $\pi$ is a surjective ring morphism.
\end{proof}

\begin{definition}\label{def:b3:rings:primemaximal}
Let $I \subsetneq A$ be a proper ideal. $I$ is
\emph{prime}\index{prime ideal} if $ab \in I \Rightarrow a \in I$
or $b \in I$; $I$ is \emph{maximal}\index{maximal ideal} if no
ideal lies strictly between $I$ and $A$.
\end{definition}

\begin{proposition}\label{prop:b3:rings:primemaximal}
$I$ is prime $\iff$ $A/I$ is an integral domain; $I$ is maximal
$\iff$ $A/I$ is a field. In particular maximal ideals are prime.
\end{proposition}

\begin{proof}
Write $\bar a$ for classes in $A/I$. ``$I$ prime'' translates
verbatim to ``$\bar a\bar b = 0 \Rightarrow \bar a = 0$ or $\bar b
= 0$'', and $A/I \neq 0$ to $I \neq A$: that is the definition of
a domain. For maximality, use the correspondence theorem: no ideal
strictly between $I$ and $A$ $\iff$ $A/I$ has no ideal other than
$0$ and itself $\iff$ $A/I$ is a field --- for the last step: in a
field the only ideals are $0$ and everything (an ideal containing
$x \ne 0$ contains $x^{-1}x = 1$); conversely if every nonzero $x$
generates the unit ideal, then $xy = 1$ for some $y$. Fields are
domains, so maximal ideals are prime.
\end{proof}

\begin{example}\label{ex:b3:rings:primemaximal}
In $\Z$: the prime ideals are $(0)$ and the $(p)$, $p$ prime; the
maximal ones are the $(p)$ ($\Z/p\Z = \mathbb F_p$ is a field,
$\Z/(0) = \Z$ is not). In $K[X, Y]$: $(X) \subsetneq (X, Y)$ are
both prime ($K[X,Y]/(X) \cong K[Y]$, a domain;
$K[X,Y]/(X,Y) \cong K$, a field), so $(X)$ is prime but not
maximal.
\end{example}

To guarantee that maximal ideals \emph{exist} in full generality,
we need a set-theoretic principle. A partially ordered set is
\emph{inductive} if every totally ordered subset (\emph{chain})
has an upper bound.

\begin{theorem}[Zorn's lemma]\label{thm:b3:rings:zorn}
Every nonempty inductive partially ordered set has a maximal
element.\index{Zorn's lemma}
\admitted
\end{theorem}

\begin{remark}\label{rem:b3:rings:zorn}
This is not a theorem of ordinary mathematics but an \emph{axiom}:
it is equivalent, over the basic Zermelo--Fraenkel axioms of set
theory, to the axiom of choice (``every product of nonempty sets
is nonempty''), which we accept throughout this book. We flag each
use. Analysis will invoke it again (Hahn--Banach,
\cref{ch:b3:banach}).
\end{remark}

\begin{theorem}[Krull]\label{thm:b3:rings:krull}
Every proper ideal $I \subsetneq A$ is contained in a maximal
ideal.\index{Krull's theorem}
\end{theorem}

\begin{proof}
Order by inclusion the set $\mathcal E$ of proper ideals
containing $I$; it is nonempty ($I \in \mathcal E$). A chain
$(J_\lambda)$ in $\mathcal E$ has upper bound $J = \bigcup
J_\lambda$: an ideal (any $a, b \in J$ lie in a common $J_\lambda$
by totality), proper ($1 \notin J_\lambda$ for all $\lambda$),
containing $I$. Zorn's lemma yields a maximal element of $\mathcal
E$, which is a maximal ideal containing $I$ (an ideal strictly
above it and proper would lie in $\mathcal E$).
\end{proof}

\begin{theorem}[Chinese remainder theorem]\label{thm:b3:rings:crt}
Let $I_1, \dots, I_n$ be pairwise \emph{comaximal}\index{comaximal
ideals} ideals of $A$ ($I_k + I_l = A$ for $k \neq l$). Then
\[
A\Big/\bigcap_{k=1}^n I_k \;\xrightarrow{\;\sim\;}\;
\prod_{k=1}^n A/I_k,
\qquad
a \longmapsto (a + I_1, \dots, a + I_n),
\]
and moreover $\bigcap_k I_k = I_1 I_2 \cdots I_n$ (the ideal
generated by products).\index{Chinese remainder theorem}
\end{theorem}

\begin{proof}
The map $f(a) = (a + I_k)_k$ is a ring morphism with kernel
$\bigcap I_k$; by \cref{thm:b3:rings:firstiso} it suffices to
prove surjectivity. Fix $k$; for each $l \neq k$ write $1 = u_l +
v_l$ with $u_l \in I_k$, $v_l \in I_l$ (comaximality). Then
\[
e_k = \prod_{l \neq k} v_l = \prod_{l\neq k}(1 - u_l)
\equiv 1 \pmod{I_k},
\qquad e_k \in I_l \ (l \neq k),
\]
so $f(e_k) = (0, \dots, 1, \dots, 0)$; given a target $(a_k +
I_k)_k$, the element $\sum_k a_k e_k$ maps to it.

Products vs intersection: $I_1\cdots I_n \subseteq \bigcap I_k$
always. Conversely, by induction it suffices to treat $n = 2$
(one checks $I_1$ and $I_2\cdots I_n$ are comaximal: multiplying
$1 = u_l + v_l$ over $l \geq 2$ gives $1 \in I_1 + I_2\cdots
I_n$). For $n = 2$: write $1 = u + v$, $u \in I_1$, $v \in I_2$;
for $x \in I_1 \cap I_2$, $x = xu + xv \in I_2I_1 + I_1I_2 =
I_1I_2$.
\end{proof}

\begin{example}\label{ex:b3:rings:crtz}
In $\Z$ with $I_k = (m_k)$, $m_k$ pairwise coprime: $\Z/(m_1\cdots
m_n)\Z \cong \prod \Z/m_k\Z$ --- the Year 2 volume's Chinese
remainder theorem. Restricting to units: $(\Z/mn\Z)^\times \cong
(\Z/m\Z)^\times \times (\Z/n\Z)^\times$ for $\gcd(m,n)=1$, whence
the multiplicativity of Euler's $\varphi$
(\cref{exo:b3:rings:8}).
\end{example}

\section{Divisibility: Euclidean, principal, factorial}

\begin{definition}\label{def:b3:rings:divisibility}
Let $A$ be an integral domain, $a, b \in A$. We say $a$
\emph{divides} $b$ ($a \mid b$) if $b \in (a) = aA$. Elements $a,
b$ are \emph{associates}\index{associate elements} if $a = ub$
with $u \in A^\times$ (equivalently $(a) = (b)$). A nonzero
nonunit $p$ is:
\begin{itemize}
  \item \emph{irreducible}\index{irreducible element} if $p = ab$
        forces $a \in A^\times$ or $b \in A^\times$;
  \item \emph{prime}\index{prime element} if $p \mid ab$ forces $p
        \mid a$ or $p \mid b$ (i.e.\ the ideal $(p)$ is prime).
\end{itemize}
\end{definition}

\begin{proposition}\label{prop:b3:rings:primeirred}
In any domain, prime $\Rightarrow$ irreducible. The converse is
false in general: in $\Z[\iu\sqrt 5] = \{a + \iu b\sqrt5 : a, b
\in \Z\}$, the element $2$ is irreducible but not prime.
\end{proposition}

\begin{proof}
Let $p$ be prime and $p = ab$. Then $p \mid ab$, say $p \mid a$:
$a = pc$, so $p = pcb$, and cancelling $p$ (domain!) gives $cb =
1$: $b \in A^\times$.

In $\Z[\iu\sqrt5]$, use the norm $N(x + \iu y\sqrt 5) = x^2 +
5y^2$, which is multiplicative (it is $\abs z^2$). If $2 = ab$
with $a, b$ nonunits, then $4 = N(a)N(b)$ with $N(a), N(b) \neq
1$ (norm-$1$ elements are $\pm1$, the units), so $N(a) = 2$:
impossible, $x^2 + 5y^2 = 2$ has no integer solution. So $2$ is
irreducible. But $2 \mid 6 = (1 + \iu\sqrt5)(1 - \iu\sqrt5)$
while $2$ divides neither factor ($\frac12 \pm \frac{\iu\sqrt5}2
\notin \Z[\iu\sqrt5]$): not prime.
\end{proof}

\begin{definition}\label{def:b3:rings:pidufd}
An integral domain $A$ is:
\begin{itemize}
  \item \emph{Euclidean}\index{Euclidean domain} if there is a map
        $\nu \colon A \setminus \{0\} \to \N$ (a \emph{Euclidean
        function}) such that for all $a, b$ with $b \ne 0$ there
        exist $q, r$ with $a = bq + r$ and ($r = 0$ or $\nu(r) <
        \nu(b)$);
  \item \emph{principal} (a \emph{PID}\index{principal ideal
        domain}) if every ideal is of the form $(a)$;
  \item \emph{factorial} (a \emph{UFD}\index{unique factorization
        domain}) if every nonzero nonunit is a product of
        irreducibles, uniquely up to order and associates.
\end{itemize}
\end{definition}

\begin{theorem}\label{thm:b3:rings:euclideanpid}
Euclidean $\Rightarrow$ principal.
\end{theorem}

\begin{proof}
Let $I \neq (0)$ be an ideal and $b \in I \setminus\{0\}$ with
$\nu(b)$ minimal. For $a \in I$, divide: $a = bq + r$; then $r = a
- bq \in I$, and $\nu(r) < \nu(b)$ would contradict minimality, so
$r = 0$ and $a \in (b)$: $I = (b)$.
\end{proof}

\begin{example}\label{ex:b3:rings:euclidean}
$\Z$ (with $\nu = \abs\cdot$) and $K[X]$ (with $\nu = \deg$) are
Euclidean --- the Year 2 volume proved both divisions. So is
$\Z[\iu]$, with $\nu = N$ the square norm
(\cref{exo:b3:rings:4}); the geometry of the proof is in the
figure below. A PID that is not Euclidean exists but is delicate
to certify (the standard example is
$\Z\bigl[\frac{1+\iu\sqrt{19}}2\bigr]$); a UFD that is not a PID
is easy: $K[X, Y]$ (\cref{exo:b3:rings:6}), or $\Z[X]$.
\end{example}

\begin{omfigure}
\begin{tikzpicture}[scale=1.05]
  \clip (-2.7,-1.7) rectangle (2.7,2.2);
  \foreach \x in {-2,...,2}
    \foreach \y in {-1,...,2}
      \fill[black!55] (\x,\y) circle (1.4pt);
  \draw[omDef, thick] (0.6,0.4) circle (1);
  \fill[omThm] (0.6,0.4) circle (1.8pt)
    node[below right, font=\small, omThm] {$a/b$};
  \fill[omDef] (1,0) circle (2pt)
    node[below, font=\small, omDef] {$q$};
  \draw[omDef, ->, shorten >=2.5pt, shorten <=2.5pt]
    (0.6,0.4) -- (1,0);
\end{tikzpicture}

{\small Division in the Gaussian integers: the exact quotient
$a/b \in \C$ lies within distance $\leq \frac{\sqrt2}{2} < 1$ of
some lattice point $q \in \Z[\iu]$; then $r = a - bq$ has $N(r) =
N(b)\,\abs{a/b - q}^2 < N(b)$. One Euclidean division, hence a
whole arithmetic.}
\end{omfigure}

\begin{lemma}[Ascending chains of principal ideals]
\label{lem:b3:rings:acc}
In a PID, every increasing sequence of ideals $I_1 \subseteq I_2
\subseteq \cdots$ is eventually constant.
\end{lemma}

\begin{proof}
$I = \bigcup_n I_n$ is an ideal (the union is increasing), so $I =
(a)$; the element $a$ lies in some $I_N$, and then $I = (a)
\subseteq I_N \subseteq I_n \subseteq I$ for $n \geq N$.
\end{proof}

\begin{lemma}[B\'ezout; Euclid's lemma]\label{lem:b3:rings:bezout}
Let $A$ be a PID and $a, b \in A$. Then $(a) + (b) = (d)$ for some
$d$, a \emph{greatest common divisor}\index{greatest common
divisor}: $d \mid a$, $d \mid b$, and every common divisor of $a,
b$ divides $d$; moreover $d = au + bv$ for some $u, v$
(B\'ezout). Consequently every irreducible element of a PID is
prime.
\end{lemma}

\begin{proof}
$(a) + (b)$ is an ideal, hence $(d)$; $a, b \in (d)$ gives $d \mid
a, b$; and $d = au + bv \in (a) + (b)$. A common divisor $c$ of
$a, b$ divides $au + bv = d$.

Euclid: let $p$ be irreducible, $p \mid ab$, $p \nmid a$. A gcd
$d$ of $p$ and $a$ divides $p$, so $d$ is a unit or an associate
of $p$ (irreducibility); associate is excluded by $p \nmid a$. So
$1 = pu + av$, whence $b = pub + abv$, and $p$ divides both terms:
$p \mid b$.
\end{proof}

\begin{theorem}\label{thm:b3:rings:pidufd}
Principal $\Rightarrow$ factorial.
\end{theorem}

\begin{proof}
\emph{Existence.} Suppose some nonzero nonunit $a$ has no
factorization into irreducibles. Then $a$ is not irreducible: $a =
a_1b_1$ with both factors nonunits; at least one of them, say
$a_1$, again has no factorization (a product of two factorizable
elements is factorizable). Iterating, we get $a = a_0, a_1, a_2,
\dots$, each a proper divisor of the last with no factorization,
so $(a_0) \subsetneq (a_1) \subsetneq (a_2) \subsetneq \cdots$
--- the inclusions are strict because $a_n = a_{n+1}c$ with $c$ a
nonunit means $(a_n) = (a_{n+1})$ would force $c \in A^\times$
(cancel in a domain). This contradicts \cref{lem:b3:rings:acc}.

\emph{Uniqueness.} Let $p_1 \cdots p_r = q_1 \cdots q_s$ with all
factors irreducible, $r \leq s$, by induction on $r$. The prime
(\cref{lem:b3:rings:bezout}) $p_1$ divides the right-hand side, so
divides some $q_j$; renumber $j = 1$. As $q_1$ is irreducible and
$p_1$ is not a unit, $q_1 = u p_1$ with $u \in A^\times$: $p_1,
q_1$ are associates. Cancel $p_1$: $p_2 \cdots p_r = (u q_2)
q_3\cdots q_s$ and conclude by induction ($r = 1$ forces $s = 1$:
a unit times irreducibles cannot be $1$).
\end{proof}

\begin{remark}\label{rem:b3:rings:ufdgcd}
In a UFD, gcds exist (take minimal exponents in the
factorizations) and Euclid's lemma holds --- irreducible $=$ prime
(\cref{exo:b3:rings:2}) --- but B\'ezout may fail: in $\Z[X]$,
$\gcd(2, X) = 1$ yet $1 \neq 2U + XV$ (evaluate at $X = 0$:
$1 = 2U(0)$, impossible). B\'ezout identities are the exclusive
property of PIDs.
\end{remark}

\begin{method}\label{met:b3:rings:quotient}
To identify a quotient ring $A/I$, hunt for a surjective morphism
$f \colon A \to B$ with kernel $I$ and invoke
\cref{thm:b3:rings:firstiso}; when $A = C[X]$ is a polynomial
ring, $f$ is usually an evaluation. Thus $\Z[X]/(X^2+1) \cong
\Z[\iu]$ (evaluate at $\iu$), $K[X,Y]/(Y - X^2) \cong K[X]$
(evaluate $Y$ at $X^2$), $\R[X]/(X^2+1) \cong \C$. To show $I$
prime or maximal, show the quotient is a domain or a field
(\cref{prop:b3:rings:primemaximal}).
\end{method}

\section{Polynomials over a UFD: Gauss and Eisenstein}

Throughout this section $A$ is a UFD with fraction field $K$
(constructed as the field of formal quotients $a/b$, $b \neq 0$,
exactly like $\Q$ from $\Z$; the Year 2 volume did this
construction for $\Q$, and it transfers verbatim). Our goal:
factoriality passes from $A$ to $A[X]$, and irreducibility over
$A$ is essentially irreducibility over the bigger field $K$.

\begin{definition}\label{def:b3:rings:content}
The \emph{content}\index{content of a polynomial} $c(P)$ of a
nonzero $P \in A[X]$ is a gcd of its coefficients (defined up to a
unit); $P$ is \emph{primitive} if $c(P) \in A^\times$. Every $P
\in A[X]$ writes $P = c(P)\,P_1$ with $P_1$ primitive, and every
$P \in K[X]\setminus\{0\}$ writes $P = \lambda P_1$ with $\lambda
\in K^\times$ and $P_1 \in A[X]$ primitive (clear denominators,
then factor out the content).
\end{definition}

\begin{lemma}[Gauss]\label{lem:b3:rings:gauss}
The product of two primitive polynomials of $A[X]$ is primitive;
consequently $c(PQ) = c(P)c(Q)$ up to units.\index{Gauss's lemma}
\end{lemma}

\begin{proof}
Let $P, Q$ be primitive and suppose some irreducible (= prime,
UFD) $p$ divides all coefficients of $PQ$. Reduce modulo $p$: in
$(A/(p))[X]$, $\bar P \bar Q = 0$. But $A/(p)$ is a domain
($(p)$ prime), so $(A/(p))[X]$ is a domain (leading coefficients
multiply), forcing $\bar P = 0$ or $\bar Q = 0$: $p$ divides all
coefficients of $P$ or all of $Q$, contradicting primitivity. For
the consequence, write $P = c(P)P_1$, $Q = c(Q)Q_1$: $PQ =
c(P)c(Q) P_1Q_1$ with $P_1Q_1$ primitive.
\end{proof}

\begin{theorem}\label{thm:b3:rings:gaussufd}
Let $A$ be a UFD with fraction field $K$.
\begin{enumerate}
  \item A primitive $P \in A[X]$ of degree $\geq 1$ is irreducible
        in $A[X]$ iff it is irreducible in $K[X]$.
  \item $A[X]$ is a UFD; its irreducibles are the irreducibles of
        $A$ and the primitive polynomials irreducible over $K$. In
        particular $\Z[X]$, and by induction $K[X_1, \dots, X_n]$
        and $\Z[X_1, \dots, X_n]$, are UFDs.
\end{enumerate}
\end{theorem}

\begin{proof}
(1) ($\Leftarrow$) If $P = QR$ in $A[X]$ with $Q, R$ nonunits,
then neither factor is constant (a constant factor of a primitive
polynomial is a unit), so the factorization is proper in $K[X]$.
($\Rightarrow$) Suppose $P = QR$ with $Q, R \in K[X]$ of degrees
$\geq 1$. Write $Q = \lambda Q_1$, $R = \mu R_1$ with $Q_1, R_1
\in A[X]$ primitive: $P = \lambda\mu\, Q_1R_1$, and $Q_1R_1$ is
primitive by Gauss. Taking contents, $\lambda\mu \in A^\times$
(both sides have unit content; formally, $\lambda\mu = c(P) \in
A^\times$ up to units, and in particular $\lambda \mu \in A$): $P
= (\lambda\mu Q_1) R_1$ is a proper factorization in $A[X]$.

(2) Existence: given $P \neq 0$ nonunit, factor $P = c(P)P_1$,
factor $c(P)$ into irreducibles of $A$, and factor $P_1$ in the
UFD $K[X]$ as $\prod Q_i$ with $Q_i \in K[X]$ irreducible; writing
$Q_i = \lambda_i R_i$ with $R_i \in A[X]$ primitive (hence
irreducible over $K$, hence in $A[X]$ by (1)), the product
$\prod \lambda_i$ is a unit of $A$ as before, and $P_1 = u\prod
R_i$. Uniqueness: compare a factorization's constant part and
polynomial part; the constants multiply to $c(P)$ (Gauss), unique
by factoriality of $A$; the polynomial parts give two
factorizations in $K[X]$ of the same polynomial, so they match up
to constants of $K^\times$ (factoriality of $K[X]$,
\cref{thm:b3:rings:pidufd}), and matching primitive polynomials
associated in $K[X]$ are associated in $A[X]$: if $R = \lambda R'$
with $R, R'$ primitive and $\lambda \in K^\times$, then taking
contents forces $\lambda \in A^\times$.
\end{proof}

\begin{theorem}[Irreducibility criteria]\label{thm:b3:rings:criteria}
Let $A$ be a UFD, $K$ its fraction field, and $P = a_nX^n + \dots
+ a_0 \in A[X]$ primitive of degree $n \geq 1$.
\begin{enumerate}
  \item (\emph{Reduction}) If $p \in A$ is prime, $p \nmid a_n$,
        and the reduction $\bar P$ is irreducible in
        $(A/(p))[X]$, then $P$ is irreducible in $K[X]$ (hence in
        $A[X]$).
  \item (\emph{Eisenstein}\index{Eisenstein's criterion}) If some
        prime $p$ satisfies $p \nmid a_n$, $p \mid a_i$ for $0
        \leq i < n$, and $p^2 \nmid a_0$, then $P$ is irreducible
        in $K[X]$ (hence in $A[X]$).
\end{enumerate}
\end{theorem}

\begin{proof}
By \cref{thm:b3:rings:gaussufd}(1), a proper factorization over
$K$ yields $P = QR$ with $Q, R \in A[X]$, $\deg Q, \deg R \geq 1$
(constants are excluded: they would be units or spoil
primitivity).

(1) Reduce mod $p$: $\bar P = \bar Q\bar R$ in $(A/(p))[X]$. Since
$p \nmid a_n$ and $\deg$ can only drop under reduction, $\deg \bar
Q = \deg Q \geq 1$ and $\deg\bar R = \deg R \geq 1$ (their leading
coefficients multiply to $\bar a_n \neq 0$, so neither drops):
$\bar P$ factors properly --- contradiction.

(2) Reduce mod $p$: $\bar Q \bar R = \bar P = \bar a_n X^n$ (all
lower coefficients die). In the domain $(A/(p))[X]$, the
factorizations of $cX^n$ ($c \ne 0$) are into constants and pure
powers $c'X^k$: indeed if $\bar Q\bar R = \bar a_nX^n$, and say
$\bar Q$ had a nonzero coefficient in degree $< \deg\bar Q$, take
lowest nonzero terms: $\operatorname{val}(\bar Q\bar R) =
\operatorname{val}\bar Q + \operatorname{val}\bar R$ (domain),
which must equal $n = \deg\bar Q + \deg\bar R$, forcing
$\operatorname{val} = \deg$ for both: both are monomials. As
above, degrees do not drop, so $Q$ and $R$ have their constant
terms $Q(0), R(0)$ divisible by $p$ --- both, since both
reductions are monomials of degree $\geq 1$. Then $p^2 \mid
Q(0)R(0) = a_0$: contradiction.
\end{proof}

\begin{example}\label{ex:b3:rings:eisenstein}
$X^n - p$ is irreducible over $\Q$ for every prime $p$ and $n
\geq 1$ (Eisenstein at $p$): there are irreducible polynomials of
every degree over $\Q$ --- in stark contrast with $\C$ (degree
$1$, d'Alembert--Gauss, proved in \cref{ch:b3:holomorphic}) and
$\R$ (degrees $1, 2$). The trick of \emph{shifting} enlarges
Eisenstein's reach: the $p$-th \emph{cyclotomic
polynomial}\index{cyclotomic polynomial} $\Phi_p = X^{p-1} +
\dots + X + 1 = \frac{X^p - 1}{X - 1}$ has
\[
\Phi_p(X + 1) = \frac{(X+1)^p - 1}{X}
= X^{p-1} + \binom{p}{1}X^{p-2} + \dots + \binom{p}{p-1},
\]
Eisenstein at $p$ ($p \mid \binom pk$ for $0 < k < p$, and
$\binom{p}{p-1} = p \not\equiv 0 \bmod p^2$): $\Phi_p(X+1)$, hence
$\Phi_p$, is irreducible over $\Q$. This is the algebraic heart of
the $17$-gon story told in \cref{ch:b3:galois}.
\end{example}

\begin{method}\label{met:b3:rings:irreducibility}
To prove $P \in \Z[X]$ irreducible over $\Q$: (i) make $P$
primitive; (ii) try Eisenstein, on $P(X)$ and on shifts $P(X \pm
1)$; (iii) try reduction modulo small primes not dividing the
leading coefficient --- irreducibility mod \emph{one} $p$
suffices, and over $\mathbb F_p$ irreducibility is a finite check
(no roots excludes degree-$1$ factors; then test the finitely many
factors of each degree $\leq \deg P/2$); (iv) if all else fails,
undetermined coefficients. Beware: reducibility mod every $p$ does
\emph{not} imply reducibility over $\Q$
(\cref{exo:b3:rings:11}).
\end{method}

\section{Noetherian rings}

\begin{definition}\label{def:b3:rings:noetherian}
A ring $A$ is \emph{Noetherian}\index{Noetherian ring} if every
ideal of $A$ is finitely generated.
\end{definition}

\begin{proposition}\label{prop:b3:rings:noethacc}
$A$ is Noetherian iff every increasing sequence of ideals is
eventually constant (\emph{ascending chain condition}), iff every
nonempty family of ideals has a maximal element (for inclusion).
\end{proposition}

\begin{proof}
\emph{(FG $\Rightarrow$ ACC)}: for a chain $I_1 \subseteq I_2
\subseteq \cdots$, the union $I$ is an ideal, generated by $x_1,
\dots, x_r$; all $x_i$ lie in some $I_N$, so $I = I_N = I_n$ for
$n \geq N$. \emph{(ACC $\Rightarrow$ maximal elements)}: if a
nonempty family $\mathcal F$ had no maximal element, pick $I_1 \in
\mathcal F$, then inductively $I_{n+1} \supsetneq I_n$ in
$\mathcal F$ (possible since $I_n$ is not maximal): an infinite
strictly increasing chain. (This uses the axiom of dependent
choices, a weak form of choice we do not fuss over.)
\emph{(Maximal elements $\Rightarrow$ FG)}: given an ideal $I$,
the family of finitely generated ideals contained in $I$ is
nonempty ($(0)$); a maximal element $J = (x_1, \dots, x_r)$ must
equal $I$: otherwise, adding $x \in I \setminus J$ to the
generators produces a strictly bigger member of the family.
\end{proof}

\begin{theorem}[Hilbert's basis theorem]\label{thm:b3:rings:hilbert}
If $A$ is Noetherian, so is $A[X]$. Hence so are $A[X_1, \dots,
X_n]$, and every quotient of them.\index{Hilbert's basis theorem}
\end{theorem}

\begin{proof}
Let $I$ be an ideal of $A[X]$, and suppose $I$ is not finitely
generated. Build a sequence: $f_1 \in I \setminus \{0\}$ of
minimal degree, and inductively $f_{k+1} \in I \setminus (f_1,
\dots, f_k)$ of minimal degree (the set is nonempty by
assumption). Degrees $d_k = \deg f_k$ are nondecreasing (by
minimality of each choice: $f_{k+1}$ was available at step
$k+1$... precisely, $f_{k+1} \notin (f_1,\dots,f_k) \supseteq
(f_1, \dots, f_{k-1})$, so $f_{k+1}$ competed at step $k$ and lost
or tied: $d_{k+1} \geq d_k$). Let $a_k \in A$ be the leading
coefficient of $f_k$. The chain of ideals $(a_1) \subseteq (a_1,
a_2) \subseteq \cdots$ stabilizes: $a_{n+1} \in (a_1, \dots,
a_n)$ for some $n$, say $a_{n+1} = \sum_{k\leq n} u_k a_k$.
Consider
\[
g = f_{n+1} - \sum_{k=1}^{n} u_k X^{\,d_{n+1} - d_k} f_k .
\]
Then $g \in I \setminus (f_1, \dots, f_n)$ (the sum lies in the
ideal, $f_{n+1}$ does not), yet the coefficient of degree
$d_{n+1}$ cancels: $\deg g < d_{n+1}$, contradicting the
minimality of $d_{n+1} = \deg f_{n+1}$.

Iterating, $A[X_1, \dots, X_n] = (A[X_1, \dots, X_{n-1}])[X_n]$
is Noetherian; a quotient $A/I$ is Noetherian because its ideals
$J/I$ lift to ideals of $A$ (correspondence), where finitely many
generators project onto generators.
\end{proof}

\begin{remark}\label{rem:b3:rings:noethmeaning}
Noetherianity is the finiteness axiom of algebraic geometry: any
system of polynomial equations in $n$ variables, however infinite,
is equivalent to finitely many of them --- its solution set is cut
out by finitely many polynomials. PIDs are Noetherian
(trivially); $\Z[X_1, X_2, \dots]$ in infinitely many variables
is not ($(X_1) \subsetneq (X_1, X_2) \subsetneq \cdots$).
Non-Noetherian rings also occur naturally in analysis: continuous
functions on $\intcc01$ form one (\cref{exo:b3:rings:10}).
\end{remark}

\section{Exercises}

\begin{exercise}[$\star$]\label{exo:b3:rings:1}
Identify the quotients: (a) $\Z[X]/(X^2 + 1) \cong \Z[\iu]$;
(b) $\R[X]/(X^2+1) \cong \C$; (c) $\mathbb F_2[X]/(X^2 + X + 1)$
is a field with $4$ elements --- write its multiplication table.
\end{exercise}

\begin{exercise}[$\star$]\label{exo:b3:rings:2}
(a) Show that in a UFD, every irreducible element is prime.
(b) Show that a finite integral domain is a field.
(c) Deduce that in a finite ring, every prime ideal is maximal.
\end{exercise}

\begin{exercise}[$\star$]\label{exo:b3:rings:3}
In $\Z[\iu\sqrt5]$: check that $3$, $1 + \iu\sqrt5$ and $1 -
\iu\sqrt5$ are irreducible, that $9 = 3\cdot 3 = (2 +
\iu\sqrt5)(2 - \iu\sqrt5)$, and conclude again (after
\cref{prop:b3:rings:primeirred}) that $\Z[\iu\sqrt5]$ is not a
UFD. Where exactly does uniqueness fail?
\end{exercise}

\begin{exercise}[$\star\star$]\label{exo:b3:rings:4}
(a) Show that $\Z[\iu]$ is Euclidean for the norm $N(x + \iu y) =
x^2 + y^2$: given $a, b \neq 0$, choose $q \in \Z[\iu]$ nearest
to $a/b \in \C$.
(b) Determine $\Z[\iu]^\times$.
(c) Same questions for $\Z[\iu\sqrt2]$ and $N(x + \iu y\sqrt2) =
x^2 + 2y^2$. Why does the same argument fail for
$\Z[\iu\sqrt5]$?
\end{exercise}

\begin{exercise}[$\star\star$]\label{exo:b3:rings:5}
Let $A$ be a ring. (a) Show that if $x$ is nilpotent ($x^n = 0$
for some $n$) then $1 + x \in A^\times$. (b) Show that if $A$ is
a domain, $A[X]^\times = A^\times$; give a counterexample over
$\Z/4\Z$. (c) Show that a domain has no idempotents ($e^2 = e$)
other than $0, 1$, and none nilpotent other than $0$.
\end{exercise}

\begin{exercise}[$\star\star$]\label{exo:b3:rings:6}
In $A = K[X, Y]$: (a) show that the ideal $(X, Y)$ is maximal but
not principal --- so $K[X,Y]$ is a UFD
(\cref{thm:b3:rings:gaussufd}) that is not a PID; (b) identify
$K[X, Y]/(Y - X^2)$ and $K[X,Y]/(XY - 1)$ as subrings of rational
functions; (c) is $(Y - X^2)$ prime? maximal?
\end{exercise}

\begin{exercise}[$\star\star$]\label{exo:b3:rings:7}
Irreducible or not over $\Q$: $X^5 - 12X^3 + 36X - 12$;\quad
$X^4 + X + 1$ \emph{(reduce mod $2$)};\quad $X^4 + 4$;\quad
$\Phi_8 = X^4 + 1$ \emph{(shift by $1$)};\quad $X^3 - X - 1$.
\end{exercise}

\begin{exercise}[$\star\star$]\label{exo:b3:rings:8}
(a) From \cref{thm:b3:rings:crt}, prove that Euler's function is
multiplicative on coprime arguments and that $\varphi(p^k) =
p^{k-1}(p-1)$; recover $\varphi(n) = n\prod_{p \mid n}(1 -
\frac1p)$.
(b) Solve: $x \equiv 2 \pmod 7$, $x \equiv 5 \pmod{11}$, $x
\equiv 1 \pmod{13}$, exhibiting the idempotents $e_k$ of the
proof of \cref{thm:b3:rings:crt}.
\end{exercise}

\begin{exercise}[$\star\star\star$]\label{exo:b3:rings:9}
(The nilradical) Let $\operatorname{Nil}(A)$ be the set of
nilpotent elements. (a) Show that $\operatorname{Nil}(A)$ is an
ideal contained in every prime ideal. (b) Conversely, let $a$ be
non-nilpotent; using Zorn's lemma on the ideals avoiding $S =
\{a^n : n \in \N\}$, produce a prime ideal not containing $a$.
Conclude:
\[
\operatorname{Nil}(A) = \bigcap_{\mathfrak p \text{ prime}}
\mathfrak p .
\]
\end{exercise}

\begin{exercise}[$\star\star\star$]\label{exo:b3:rings:10}
(a) Let $A$ be Noetherian and $f \colon A \to A$ a surjective
ring morphism. Show that $f$ is injective. \emph{(Consider $\ker
f \subseteq \ker f^2 \subseteq \cdots$.)}
(b) Show that the ring $\mathcal C(\intcc01, \R)$ of continuous
functions is not Noetherian. \emph{(Consider $I_n = \{f : f = 0
\text{ on } \intcc0{1/n}\}$.)}
(c) Show that in a Noetherian \emph{domain}, every nonzero
nonunit is a (finite) product of irreducibles --- so
non-factoriality of $\Z[\iu\sqrt 5]$ is a failure of
\emph{uniqueness} only.
\end{exercise}

\begin{exercise}[$\star\star\star$]\label{exo:b3:rings:11}
Let $P = X^4 + 1$.
(a) Show that $P$ is irreducible over $\Q$
(\cref{exo:b3:rings:7}).
(b) Show that $P$ is reducible modulo \emph{every} prime $p$:
treat $p = 2$; then, for odd $p$, show that $8 \mid p^2 - 1$ and
admit for now (proved in \cref{ch:b3:galois}) that the
multiplicative group of the field with $p^2$ elements is cyclic,
to conclude that $P$ splits into two quadratic factors mod $p$;
make them explicit when one of $-1$, $2$, $-2$ is a square mod
$p$, and show one of them always is.
\end{exercise}

\section{Problem: Fermat's two-squares theorem}

\begin{problem}[{Weekend problem --- sums of two squares, via
$\Z[\iu]$}]\label{pb:b3:rings:1}
Which integers are sums of two squares? Fermat's answer (1640) is
one of arithmetic's gems; the Gaussian integers turn its proof
into ring theory. Throughout, $N(x + \iu y) = x^2 + y^2$ denotes
the norm, $\Z[\iu]$ is Euclidean (\cref{exo:b3:rings:4}), hence a
PID and a UFD, and \emph{Gaussian prime} means prime (=
irreducible) element of $\Z[\iu]$.

\medskip
\noindent\textbf{Part I --- Norms and Gaussian primes.}
\begin{enumerate}
  \item Verify $N(zw) = N(z)N(w)$, deduce again $\Z[\iu]^\times =
        \{\pm1, \pm\iu\}$, and prove the \emph{Brahmagupta
        identity}: a product of two sums of two squares is a sum
        of two squares.
  \item Show that if $N(z)$ is a prime number, then $z$ is a
        Gaussian prime.
  \item Show that every Gaussian prime $\pi$ divides exactly one
        prime number $p$ \emph{(consider $N(\pi) = \pi\bar\pi$)},
        and that then $N(\pi) \in \{p, p^2\}$.
  \item Deduce the dichotomy: for each prime $p$, either $p$
        stays prime in $\Z[\iu]$ (and no Gaussian prime of norm
        $p$ exists), or $p = \pi\bar\pi$ with $\pi$ a Gaussian
        prime of norm $p$ --- and then $p = a^2 + b^2$.
\end{enumerate}

\medskip
\noindent\textbf{Part II --- Wilson's theorem and $-1$ modulo
$p$.}
\begin{enumerate}[resume]
  \item Prove \emph{Wilson's theorem}: for $p$ prime, $(p-1)!
        \equiv -1 \pmod p$. \emph{(Pair each residue with its
        inverse; which ones are self-paired?)}
  \item Let $p$ be an odd prime and $m = \frac{p-1}2$. Show that
        $(m!)^2 \equiv (-1)^{m+1} \pmod p$ \emph{(in $(p-1)!$,
        replace each factor $k > m$ by $-(p - k)$)}.
  \item Conclude: $-1$ is a square modulo $p$ iff $p = 2$ or $p
        \equiv 1 \pmod 4$. \emph{(For the ``only if'': if $x^2
        \equiv -1$, what is the order of $x$ in $(\Z/p\Z)^\times$,
        and what does Lagrange say?)}
\end{enumerate}

\medskip
\noindent\textbf{Part III --- The splitting law.}
\begin{enumerate}[resume]
  \item Let $p \equiv 1 \pmod 4$, and $x$ with $p \mid x^2 + 1 =
        (x + \iu)(x - \iu)$. Show that $p$ is \emph{not} a
        Gaussian prime, and conclude with Part I: $p = a^2 + b^2$.
  \item Let $p \equiv 3 \pmod 4$. Show directly that $p$ is not a
        sum of two squares \emph{(squares mod $4$)}, and deduce
        that $p$ stays a Gaussian prime.
  \item Settle $p = 2$: exhibit the factorization $2 =
        -\iu(1+\iu)^2$ and check that $1 + \iu$ is a Gaussian
        prime. ($2$ is the unique \emph{ramified} prime: divisible
        by the square of a Gaussian prime up to a unit.)
  \item Assemble the classification of Gaussian primes, up to
        units: $1 + \iu$; the integers $p \equiv 3 \pmod 4$; the
        conjugate pairs $\pi, \bar\pi$ of norm $p \equiv 1 \pmod
        4$. Verify it on $5 = (2+\iu)(2-\iu)$ and on $3$.
\end{enumerate}

\medskip
\noindent\textbf{Part IV --- The two-squares theorem.}
\begin{enumerate}[resume]
  \item Prove the direct half: if in the factorization $n =
        \prod p_i^{\alpha_i}$ every prime $\equiv 3 \pmod 4$
        appears with an even exponent, then $n$ is a sum of two
        squares. \emph{(Brahmagupta + Parts II--III.)}
  \item Prove the converse: if $n = a^2 + b^2 = N(a + \iu b)$ and
        $q \equiv 3 \pmod 4$ divides $n$, show that $q$, a
        Gaussian prime, divides $a + \iu b$ or $a - \iu b$, that
        it in fact divides both $a$ and $b$, and conclude by
        induction on $n$ that the exponent of $q$ in $n$ is even.
  \item State the final theorem. Which of $2025$, $2026$, $2027$
        are sums of two squares? \emph{($2025 = 81 \cdot 25$;
        $2026 = 2 \cdot 1013$, $1013$ prime; $2027$ prime.)}
  \item (Epilogue) Show that a prime $p \equiv 1 \pmod 4$ is a
        sum of two squares in an essentially unique way: if $p =
        a^2 + b^2 = c^2 + d^2$ (positive integers), then $\{a, b\}
        = \{c, d\}$. \emph{(Uniqueness of factorization in
        $\Z[\iu]$.)}
\end{enumerate}
\end{problem}
